% =====================================================================
%  Legend of the Technomancers — MECHANICS DESIGN SPACE
%  A brainstorm document. Not rules. Not decisions.
%  Build:  tectonic -X compile design-space.tex
% =====================================================================
\documentclass[11pt,oneside]{memoir}
\usepackage{atrpg}

% Wider single-column measure than the rulebook — this document is all
% comparison tables, and they need the room.
\geometry{a4paper, inner=18mm, outer=18mm, top=20mm, bottom=22mm}

% atrpg.sty loads tikz but not this library; the dependency diagram in the
% final chapter uses relative "below=of" placement.
\usetikzlibrary{positioning}

% Option tables must break across pages — see the opttable environment.
\usepackage{longtable}

% ---- Local macros, this document only -------------------------------

% Crunch meter, 1–5 filled blocks. "How much rules weight does this add?"
\newcommand{\crunch}[1]{%
  $\vcenter{\hbox{\begin{tikzpicture}
    \foreach \i in {1,...,5}{%
      \ifnum\i>#1
        \draw[draw=atrRule,line width=0.4pt] (\i*5pt,0) rectangle ++(3.2pt,7pt);
      \else
        \fill[\realmaccent] (\i*5pt,0) rectangle ++(3.2pt,7pt);
      \fi}
  \end{tikzpicture}}}$}

% Inline source citation for an inspiring game.
\newcommand{\src}[1]{{\small\itshape\color{atrAccent2}#1}}

% Option label — the bold short name in the first column of every table.
% \textbf (not a bare \bfseries group) because \textbf supplies \leavevmode.
% A font switch in vertical mode at the head of a p-cell pushes the cell a
% whole line below its neighbours — see \thc below for the same problem.
\newcommand{\optname}[1]{\textbf{\color{\realmaccent}#1}}

% Table header cell. \hdfont is a raw font switch, so it needs \leavevmode
% explicitly or the header row sits a line low.
\newcommand{\thc}[1]{\leavevmode{\hdfont\bfseries\color{white}#1}}

% Standard four-column option table.
%
% longtable, for two reasons:
%  1. These tables run 6–9 rows and must be able to break across a page;
%     a plain tabular is unbreakable and orphans its own title.
%  2. tabularx is unusable here regardless — its body scanner looks for a
%     literal \end{tabularx} and never finds one hidden inside
%     \end{opttable}. (Same constraint as gearlist in atrpg.sty.)
% Hence fixed p-widths, which must sum to this document's \textwidth.
% The title is a table row, not a paragraph above the table, so it can
% never be separated from the header.
\newenvironment{opttable}[1]{%
  \par\smallskip
  \setlength{\tabcolsep}{5pt}\renewcommand{\arraystretch}{1.28}
  \setlength{\LTleft}{0pt}\setlength{\LTright}{\fill}
  \begin{longtable}{@{}>{\leavevmode\RaggedRight}p{3.2cm} >{\leavevmode\RaggedRight}p{8.8cm} >{\leavevmode\RaggedRight}p{2.9cm} >{\leavevmode\centering\arraybackslash}p{1.4cm}@{}}
  \multicolumn{4}{@{}l@{}}{\leavevmode{\hdfont\normalsize\bfseries\color{\realmaccent}#1}}\\[2pt]
  \rowcolor{\realmaccent}
  \thc{Option} & \thc{How it works} & \thc{Inspired by} & \thc{Weight}\\
  \endfirsthead
  \multicolumn{4}{@{}l@{}}{\leavevmode{\hdfont\normalsize\bfseries\color{\realmaccent}#1 \textnormal{\itshape(continued)}}}\\[2pt]
  \rowcolor{\realmaccent}
  \thc{Option} & \thc{How it works} & \thc{Inspired by} & \thc{Weight}\\
  \endhead}
{\end{longtable}\par\medskip}

% A row of the option table.
\newcommand{\opt}[4]{\optname{#1} & #2 & \src{#3} & \crunch{#4}\\}

% Alternating row tint for readability in long tables.
\newcommand{\tintrow}{\rowcolor{atrBoxbg}}

% Verdict-free observation box — points out interactions, doesn't pick winners.
\newtcolorbox{dread}[1][Designer's read]{%
  enhanced, breakable,
  colback=atrPaper, colframe=\realmaccent, boxrule=0pt,
  borderline west={2.2pt}{0pt}{\realmaccent},
  left=10pt, right=10pt, top=6pt, bottom=6pt,
  fonttitle=\hdfont\bfseries,
  coltitle=\realmaccent, colbacktitle=atrPaper, title={#1}}

\begin{document}
\color{atrInk}

% ------------------------------------------------------------ Title page
\begin{titlingpage}
\begin{tikzpicture}[remember picture, overlay]
  \fill[atrInk] (current page.north west) rectangle ([yshift=-118mm]current page.north east);
  \node[anchor=north, align=center, inner sep=0] at ([yshift=-30mm]current page.north){%
    \thc{\fontsize{34}{38}\selectfont LEGEND OF THE}\\[2mm]
    \thc{\fontsize{34}{38}\selectfont TECHNOMANCERS}\\[8mm]
    {\color{atrAccent2}\rule{0.50\paperwidth}{1.2pt}}\\[8mm]
    {\hdfont\bfseries\color{atrAccent2}\fontsize{30}{34}\selectfont THE DESIGN SPACE}\\[7mm]
    {\hdfont\large\itshape\color{white}Options for the Engine --- A Brainstorm, Not a Decision}%
  };
  \node[anchor=south, align=center, inner sep=0] at ([yshift=26mm]current page.south){%
    {\itshape\color{atrInk}Every system in this document is a live candidate.}\\[2pt]
    {\itshape\color{atrInk}Nothing here is a rule yet.}\\[10mm]
    {\hdfont\color{atrRule}INTERNAL DESIGN DOCUMENT}%
  };
\end{tikzpicture}
\end{titlingpage}

\frontmatter
\tableofcontents

\mainmatter

% =====================================================================
\chapter{Reading This Document}
% =====================================================================

\lettrine{T}{his is a menu}, not a proposal. Each chapter takes one
question the game has not yet answered, lays out every serious way to
answer it, and shows what each answer costs. Options are not ranked.
Where an option interacts badly --- or unexpectedly well --- with another
part of the game, that is flagged in a \emph{Designer's read}.

\section{What Is Already Locked}

These are settled and every option in this document must live with them.

\begin{statblock}[Fixed Constraints]
\begin{tabularx}{\linewidth}{@{}p{4.4cm} Y@{}}
\keyword{Four Columns} & Body, Stuff, Skills, Magic.\\[3pt]
\keyword{Ranks 1--4} & Each column is rated \rating{1} to \rating{4}.\\[3pt]
\keyword{Priority Assignment} & One of A/B/C/D to each column, no repeats. You cannot be good at everything.\\[3pt]
\keyword{Multi-Genre} & The same sheet must work in high fantasy, cyberpunk, and space opera without rewriting.\\[3pt]
\keyword{Rank Travels} & The number stays; the fiction it describes is local to the realm.\\[3pt]
\keyword{Characters Differ} & The reason Body and Stuff are pillars is that they are \emph{cool} --- they capture the ways characters can be interestingly unlike one another. Any candidate that flattens that has failed, whatever its other virtues.\\
\end{tabularx}
\end{statblock}

\section{What Is \emph{Not} Locked, Despite Appearances}

Two things the current draft does are easy to mistake for constraints.
Both are merely consequences of one early decision, and both are open.

\begin{statblock}[Open, Not Settled]
\begin{tabularx}{\linewidth}{@{}p{4.4cm} Y@{}}
\keyword{The 0--10 Pool} & Not a cap and not a target --- just what falls out of ``two traits, each 0--5.'' Change the traits and the range changes. Where this document quotes a ceiling of \dice{10}, read it as a description of the current draft, never as a requirement.\\[3pt]
\keyword{Rules-Light} & Also unsettled. A heavier game with several genuinely distinct subsystems is on the table, and may well be \emph{better} at making characters differ. Treat the weight meters below as a price tag, not a disqualification.\\
\end{tabularx}
\end{statblock}

\section{What Is Still Open}

Everything about how those numbers become dice. Specifically:

\begin{enumerate}
  \item What the \textbf{core roll} is, and whether it needs a GM-set target.
  \item How \textbf{gear} enters a roll --- the question that started this document.
  \item What each of the \textbf{four columns} actually hands you at creation.
  \item How \textbf{harm} works, and what armour does about it.
\end{enumerate}

\section{The Weight Dial}

Every option is tagged with a weight rating. This is not quality --- it is
how much rules-mass the option adds, measured roughly as ``how much must a
player or GM hold in their head to use this correctly.''

\begin{center}
\begin{tikzpicture}
  \shade[left color=columnSkills, right color=columnBody]
    (0,0) rectangle (14,0.55);
  \foreach \x/\l in {1.4/1, 4.2/2, 7/3, 9.8/4, 12.6/5}{
    \draw[white, line width=1pt] (\x,0) -- (\x,0.55);
    \node[below=2pt, font=\hdfont\small] at (\x,0) {\l};}
  \node[above=3pt, font=\hdfont\small\color{columnSkills}, anchor=west] at (0,0.55)
    {FEATHERWEIGHT};
  \node[above=3pt, font=\hdfont\small\color{columnBody}, anchor=east] at (14,0.55)
    {FULL SIMULATION};
  \node[below=22pt, font=\footnotesize\itshape, anchor=west] at (0,0)
    {A rule you never look up};
  \node[below=22pt, font=\footnotesize\itshape, anchor=east] at (14,0)
    {A rule with its own chapter};
\end{tikzpicture}
\end{center}

\begin{dread}[How to use this]
Read Chapters 2 and 3 first --- the core roll and the gear question
constrain everything downstream. Chapters 4 through 7 are largely
independent of each other and can be mixed freely. Chapter 9 assembles
six coherent whole-game packages from these parts, if picking
\'a~la~carte is not appealing.
\end{dread}

% =====================================================================
\renewcommand{\realmaccent}{atrAccent}
\chapter{The Core Roll}
% =====================================================================

\lettrine{E}{verything} downstream depends on this. The current draft
uses a d6 dice pool counting 4--6 as successes against a GM-set target.
That is one of at least eight viable engines.

\begin{opttable}{Eight Core Engines}
\tintrow\opt{Count Successes}{Pool of d6. Each 4, 5, or 6 is a success. GM names how many successes the task needs. \emph{This is the current draft.}}{Vampire: the Masquerade, Shadowrun, Exalted}{3}
\opt{Highest Die}{Roll the pool, read only the single highest die. 6 = full success, 4--5 = success with a cost, 1--3 = failure. Two or more 6s = critical.}{Blades in the Dark, Forged in the Dark}{2}
\tintrow\opt{Roll \& Keep}{Roll a number of dice equal to one trait, keep the best N equal to the other trait, sum them.}{Legend of the Five Rings, 7th Sea}{4}
\opt{Step Dice}{Every trait is a \emph{die size}, not a count: d4 through d12. Roll two traits, sum, compare to a difficulty.}{Cortex Prime, Savage Worlds}{3}
\tintrow\opt{2d6 + Modifier}{Always roll exactly 2d6 and add a small modifier. 6 or less = miss, 7--9 = partial, 10+ = full.}{Apocalypse World, Traveller}{1}
\opt{Opposed Pools}{Both sides build a pool and roll; compare success counts. No GM-set difficulty --- the world rolls back.}{Burning Wheel, Legend of the Five Rings}{4}
\tintrow\opt{Variable Target}{Pool is fixed by traits; \emph{difficulty changes the number each die must beat} (3+ easy, 5+ hard) rather than how many successes you need.}{Shadowrun 3e, original Storyteller}{3}
\opt{Roll Under}{Roll one die under Stat + Skill on a fixed scale. Simple, but no degrees of success without extra rules.}{Call of Cthulhu, Cyberpunk}{2}
\end{opttable}

\section{The Numbers Behind Option 1}

If the pool stays at \dice{0}--\dice{10} and 4--6 succeeds, each die is a
coin flip. Expected successes are exactly half the pool. Here is the
actual chance of hitting a given target:

\begin{center}
\setlength{\tabcolsep}{5pt}\renewcommand{\arraystretch}{1.25}
\begin{tabular}{@{}c|cccccc@{}}
\rowcolor{\realmaccent}
\thc{Pool} &
\thc{1+} & \thc{2+} &
\thc{3+} & \thc{4+} &
\thc{5+} & \thc{6+}\\
\dice{2}  & 75\% & 25\% & --- & --- & --- & ---\\
\rowcolor{atrBoxbg}
\dice{3}  & 88\% & 50\% & 13\% & --- & --- & ---\\
\dice{4}  & 94\% & 69\% & 31\% & 6\% & --- & ---\\
\rowcolor{atrBoxbg}
\dice{5}  & 97\% & 81\% & 50\% & 19\% & 3\% & ---\\
\dice{6}  & 98\% & 89\% & 66\% & 34\% & 11\% & 2\%\\
\rowcolor{atrBoxbg}
\dice{7}  & 99\% & 94\% & 77\% & 50\% & 23\% & 6\%\\
\dice{8}  & 99\% & 97\% & 86\% & 64\% & 36\% & 15\%\\
\rowcolor{atrBoxbg}
\dice{9}  & 99\% & 98\% & 91\% & 75\% & 50\% & 25\%\\
\dice{10} & 99\% & 99\% & 95\% & 83\% & 62\% & 38\%\\
\end{tabular}
\end{center}

Read the diagonal: a target of \emph{half your pool} is always a coin
flip. That is a genuinely elegant property --- the GM can set difficulty
by eye. Targets of 1 and 2 are nearly free above \dice{5} dice, which
means the useful difficulty band is narrow: roughly 3 to 5 successes.

\section{The Numbers Behind Option 2}

Highest-die reading needs \emph{small} pools, because a large pool almost
always contains a 6:

\begin{center}
\setlength{\tabcolsep}{6pt}\renewcommand{\arraystretch}{1.25}
\begin{tabular}{@{}c|ccc@{}}
\rowcolor{\realmaccent}
\thc{Pool} &
\thc{Full (a 6)} &
\thc{Partial (4--5)} &
\thc{Miss}\\
\dice{1} & 17\% & 33\% & 50\%\\
\rowcolor{atrBoxbg}
\dice{2} & 31\% & 44\% & 25\%\\
\dice{3} & 42\% & 45\% & 13\%\\
\rowcolor{atrBoxbg}
\dice{4} & 52\% & 42\% & 6\%\\
\dice{6} & 67\% & 32\% & 2\%\\
\rowcolor{atrBoxbg}
\dice{10} & 84\% & 16\% & 1\%\\
\end{tabular}
\end{center}

\begin{dread}
This is the central tension of the whole document. \textbf{Counting
successes scales beautifully to ten dice but requires the GM to invent a
number on every single roll.} \textbf{Highest-die never requires a
number and gives graded outcomes for free, but saturates above about
four dice.}

Adopting highest-die therefore means abandoning the \dice{0}--\dice{10}
range and rescaling the whole game to roughly \dice{0}--\dice{4}. That is
not a small change --- but it would make the four columns much easier to
balance, because there is far less room for them to drift apart.
\end{dread}

\section{Partial Success}

Worth separating out, because it is orthogonal to the engine. The current
draft has only pass/fail. Three ways to add gradation:

\begin{opttable}{Adding Degrees of Outcome}
\tintrow\opt{Margin}{Successes above the target are ``extra effect'' --- more damage, faster, quieter.}{Shadowrun, Exalted}{1}
\opt{Miss by One}{Falling exactly one success short succeeds anyway, at a cost the GM names.}{Genesys, many house rules}{1}
\tintrow\opt{Position \& Effect}{Before rolling, the GM names how dangerous the attempt is and how much it accomplishes. The roll resolves both.}{Blades in the Dark}{3}
\end{opttable}

% =====================================================================
\renewcommand{\realmaccent}{columnStuff}
\chapter{Where Gear Enters the Roll}
% =====================================================================

\lettrine{T}{his is the question} that prompted the document. The
instinct was: \emph{a driving roll could be a number for the car plus a
number for your skill.} That is one of nine answers, and the choice here
decides what the entire Stuff column is for.

To make them comparable, every option below is applied to the same
fiction: \textbf{Kira, Agility \dice{3}, Piloting \dice{2}, driving a
rank-\dice{4} vehicle through a chase.}

\begin{opttable}{Nine Ways Gear Can Matter}
\tintrow\opt{Substitution}{The gear's rating \emph{replaces} the Stat. Pool = Gear + Skill. Bad gear actively limits a good driver. \newline \textbf{Kira rolls 4 + 2 = \dice{6}.}}{Ars Magica tools, many hacks}{2}
\opt{Best-Of}{Pool = higher of (Stat, Gear) + Skill. Gear is pure upside; gear below your Stat is invisible. \newline \textbf{Kira rolls 4 + 2 = \dice{6}.}}{Numenera assets}{2}
\tintrow\opt{Additive, Capped}{Pool = Stat + Skill + Gear, total capped at \dice{10}. Simple, but pushes almost everyone to the cap. \newline \textbf{Kira rolls 3 + 2 + 4 = \dice{9}.}}{Shadowrun}{2}
\opt{Quality, Not Quantity}{Gear lowers the \emph{success threshold} instead of adding dice. Rank 3+ gear: 3--6 succeed, not 4--6. Pool size never changes. \newline \textbf{Kira rolls \dice{5}, succeeding on 3+.}}{Shadowrun 3e, Storyteller difficulty}{3}
\tintrow\opt{Automatic Successes}{Gear rank adds free successes \emph{after} the roll. Raises the floor, not the ceiling --- good gear means never fumbling. \newline \textbf{Kira rolls \dice{5}, then adds 2.}}{Exalted, Mutants \& Masterminds}{2}
\opt{Re-Rolls}{Gear rank = how many failed dice you may re-roll. Same average as adding dice, much lower variance. \newline \textbf{Kira rolls \dice{5}, re-rolls up to 4 misses.}}{Genesys, 2d20 momentum}{3}
\tintrow\opt{Gear as Tags}{Items carry descriptive tags. Applicable tags add dice, exactly as Magic already does. \emph{Unifies Stuff and Magic into one subsystem.} \newline \textbf{\emph{fast}, \emph{armoured}, \emph{familiar} apply: \dice{6}.}}{Lady Blackbird, City of Mist}{2}
\opt{Permission Only}{Gear grants no dice at all. It decides \emph{what you may attempt} and how big the effect is. No car, no car chase. \newline \textbf{Kira rolls 3 + 2 = \dice{5}, but can attempt things a pedestrian cannot.}}{Blades in the Dark, Apocalypse World}{1}
\tintrow\opt{Separate Asset Die}{Gear is its own die rolled alongside the others and added to the pool as a distinct object that can be damaged or lost. \newline \textbf{Kira rolls Agi + Piloting + a vehicle die.}}{Cortex Prime}{4}
\end{opttable}

\section{What Each Does to the Maths}

All applied to a baseline \dice{6}-dice pool, so the shapes are
comparable. Note that three different options land on the same average
and behave completely differently:

\begin{center}
\setlength{\tabcolsep}{6pt}\renewcommand{\arraystretch}{1.28}
\begin{tabularx}{\textwidth}{@{}p{4.4cm}ccc Y@{}}
\rowcolor{\realmaccent}
\thc{Gear effect} &
\thc{Avg} &
\thc{Spread} &
\thc{Floor} &
\thc{Character of the change}\\
Baseline: \dice{6} dice & 3.0 & 1.22 & 0 & ---\\
\rowcolor{atrBoxbg}
+2 dice (additive) & 4.0 & 1.41 & 0 & \textbf{Raises the ceiling.} Swingier; big wins become possible.\\
Threshold to 3+ & 4.0 & 1.15 & 0 & \textbf{Tightens the spread.} Reliable, not spectacular.\\
\rowcolor{atrBoxbg}
+2 auto successes & 5.0 & 1.22 & 2 & \textbf{Raises the floor.} You can no longer fail badly.\\
Re-roll 2 misses & 4.0 & 1.17 & 0 & \textbf{Kills the bad tail.} Feels like quality without inflation.\\
\end{tabularx}
\end{center}

\begin{dread}
``Better gear'' is not one idea --- it is at least three. A legendary
sword that \emph{raises your ceiling} plays completely differently from
one that \emph{raises your floor}, even at identical average output. It
is worth deciding what owning excellent equipment is supposed to
\emph{feel} like before picking a mechanism.

Separately: the \textbf{Gear as Tags} option is the only one that collapses
two of the four columns into a single engine. If Stuff and Magic both run
on ``two dice per applicable tag,'' the game has one rule instead of two,
and the Stuff chapter stops needing a price list at all.
\end{dread}

\section{The Substitution Family in Detail}

If substitution wins, a follow-up question appears: gear can stand in for
\emph{either} side of the roll, and which side it replaces says something
about the item.

\begin{center}
\setlength{\tabcolsep}{6pt}\renewcommand{\arraystretch}{1.3}
\begin{tabularx}{\textwidth}{@{}p{3.0cm} p{4.0cm} Y@{}}
\rowcolor{\realmaccent}
\thc{Kind of thing} &
\thc{What it replaces} &
\thc{Examples}\\
\optname{Instrument} & The \textbf{Stat}. The tool has the power; you supply the craft. & Sword, rifle, car, lockpicks, a good telescope.\\
\rowcolor{atrBoxbg}
\optname{Automaton} & The \textbf{Skill}. The tool has the craft; you supply the will. & Smart-targeting gun, autopilot, familiar, hireling, an enchanted map.\\
\optname{Protection} & Nothing --- static effect. & Armour, wards, a shield generator.\\
\rowcolor{atrBoxbg}
\optname{Holding} & Nothing --- changes what is possible. & Base, ship, fortress, wealth, a standing army.\\
\end{tabularx}
\end{center}

This taxonomy is independent of which substitution rule is chosen, and it
gives the Stuff catalogue a spine: every item is exactly one of these four.

% =====================================================================
\renewcommand{\realmaccent}{columnBody}
\chapter{Body}
% =====================================================================

\lettrine{B}{ody} currently sets four Stats (Strength, Agility, Charisma,
Intellect) rated \dice{0}--\dice{5}, with Strength doubling as durability.
Ranks \rating{3} and \rating{4} are undesigned.

\begin{opttable}{Eight Ways to Build a Body}
\tintrow\opt{Dots + Cap}{A budget of dots to spread across four Stats, plus a cap on any one. \emph{Cap = rank + 1} makes it one rule.}{Vampire: the Masquerade}{3}
\opt{Modification Packages}{No free-form points. Pick N traits from a list; each is a bundle. The current rank-\rating{2} rule, extended to all ranks.}{PbtA playbooks, Lancer}{2}
\tintrow\opt{Fixed Arrays}{Each rank hands you a fixed set of numbers to arrange as you like --- rank \rating{4} = [5,4,2,1]. Zero arithmetic.}{D\&D standard array, Fate pyramid}{1}
\opt{Forms}{Pick a \emph{Form} --- Human, Ogre, Sylph, Dragon --- that bundles stats and traits together, and re-skins per realm.}{Savage Worlds races, M\&M archetypes}{2}
\tintrow\opt{Collapse to One}{Drop the four Stats entirely. Body \emph{is} your physical number. Pairs with a skill-heavy or approach-based design.}{Risus, Lasers \& Feelings}{1}
\opt{Two or Three Stats}{Physical / Social / Mental instead of four. Fewer dump stats, faster sheets, less character texture.}{Cypher System, many indies}{2}
\tintrow\opt{Body as Die Size}{Each Stat is a die size set by Body rank: rank \rating{2} gives you a d8 and three d6s.}{Cortex Prime}{3}
\opt{Body Caps Only}{Everyone gets the same dot budget; Body rank sets only the \emph{ceiling}. High Body = specialisation, not more total power.}{Original idea}{2}
\end{opttable}

\section{The Durability Question}

Separate and easy to overlook: \textbf{Strength currently does two jobs.}
That coupling forces every character who wants to survive a fight to buy
Strength, which quietly punishes the clever and the charming.

\begin{opttable}{What Measures How Much Harm You Take}
\tintrow\opt{Strength}{The current draft. Strong characters are tough. Simple, but couples offence and survival.}{Vampire (Stamina), many}{1}
\opt{Body Rank Itself}{Body \rating{3} = three boxes of harm, whatever your Stats are. Cleanly splits the column's two jobs: \emph{rank is survival, dots are dice}.}{Original idea}{1}
\tintrow\opt{Flat + Rank}{Everyone has a base track; Body rank extends it. Keeps low-Body characters playable.}{PbtA, Fate}{1}
\opt{No Track At All}{Harm is fictional conditions, not a number. Body rank determines how many you can carry.}{Apocalypse World}{2}
\end{opttable}

\begin{dread}
\textbf{Body rank as the harm track} is the cheapest good idea in this
chapter. It costs one sentence, it uncouples toughness from Strength, and
it gives the Body column a second distinct function that no other column
can provide --- which is exactly what a four-column game needs in order to
keep all four columns worth buying.
\end{dread}

\section{Sketching Ranks 3 and 4}

Whatever mechanism is chosen, ranks \rating{3} and \rating{4} need to be
\emph{qualitatively} different, not just larger. Three directions:

\begin{opttable}{Making High Body Strange}
\tintrow\opt{More of the Same}{Higher caps and more dots. Rank \rating{4} simply means bigger numbers.}{Most point-buy games}{1}
\opt{Unlocked Traits}{Rank \rating{3} unlocks \emph{Exceptional} traits (uncanny senses, inhuman reflexes); rank \rating{4} unlocks \emph{Monstrous} ones (wings, natural weapons, size, armoured hide, not needing to breathe).}{Exalted, Monster of the Week}{2}
\tintrow\opt{Change the Rules}{High Body breaks a rule rather than raising a number --- you ignore falling damage, you do not need to eat, mundane weapons cannot hurt you.}{Nobilis, Exalted Charms}{3}
\end{opttable}

% =====================================================================
\renewcommand{\realmaccent}{columnStuff}
\chapter{Stuff}
% =====================================================================

\lettrine{S}{tuff} is the least-designed column and the one most exposed
to the multi-genre problem. Whatever system is chosen must survive a
character walking out of a dungeon and into a space station.

\begin{opttable}{Eight Ways to Own Things}
\tintrow\opt{Dots = Cost = Rating}{Rank gives a dot budget; an item's rating \emph{is} its price. A rank-3 weapon costs 3 dots and rates 3. One number per item instead of two.}{Vampire backgrounds}{3}
\opt{N Picks}{Rank \emph{N} lets you choose \emph{N} items of rating \emph{N} or lower from a catalogue. No arithmetic whatsoever.}{Lancer licences, PbtA gear lists}{1}
\tintrow\opt{Money}{Each rank grants a cash sum; items have prices. Maximum granularity, maximum bookkeeping, and cash does not travel between realms.}{Shadowrun, Traveller}{4}
\opt{Load}{You never declare gear in advance. You carry a Load rating; mid-scene you say ``of course I brought rope'' and mark it off. \emph{Gear re-skins itself automatically because it is not named until used.}}{Blades in the Dark}{2}
\tintrow\opt{Signature Items}{Name one to three items that are genuinely yours and travel with you, re-skinning per realm. Everything else is assumed.}{Fate stunts, Cortex signature assets}{1}
\opt{Gear as Tags}{Items are bundles of tags; applicable tags add dice. Identical engine to Magic.}{Lady Blackbird, City of Mist}{2}
\tintrow\opt{Wealth Rating}{You do not own a list --- you have a Wealth number and roll it to see whether you happen to have or can obtain a thing.}{Cypher System, Fate}{2}
\opt{World Claims}{Stuff does not buy objects, it buys \emph{facts about the world}: a ship, a safehouse, a contact, turf, a reputation. It changes the campaign rather than the roll.}{Blades crew sheet, Ars Magica covenant}{3}
\end{opttable}

\section{The Realm-Crossing Filter}

This is the test that matters most for Stuff. A character crosses from
high fantasy into the sprawl. What happens to their possessions?

\begin{center}
\setlength{\tabcolsep}{6pt}\renewcommand{\arraystretch}{1.3}
\begin{tabularx}{\textwidth}{@{}p{3.4cm} >{\leavevmode\centering\arraybackslash}p{2.6cm} Y@{}}
\rowcolor{\realmaccent}
\thc{System} &
\thc{Survives crossing?} &
\thc{What actually happens at the table}\\
Money & \textcolor{columnBody}{\bfseries Badly} & Gold is worthless; you must re-price everything and re-shop. Kills pace on arrival.\\
\rowcolor{atrBoxbg}
Dots = Cost & \textcolor{atrAccent2}{\bfseries Workable} & Lines on the sheet stay; you re-skin each one. Requires a translation table per realm --- which the genre chapters already do.\\
N Picks & \textcolor{atrAccent2}{\bfseries Workable} & Same as above, with less arithmetic.\\
\rowcolor{atrBoxbg}
Load & \textcolor{columnSkills}{\bfseries Cleanly} & Nothing to translate. You never wrote down ``sword'' in the first place.\\
Signature Items & \textcolor{columnSkills}{\bfseries Cleanly} & Your one named thing becomes its local equivalent. This is arguably the \emph{theme} of the game.\\
\rowcolor{atrBoxbg}
Gear as Tags & \textcolor{columnSkills}{\bfseries Cleanly} & \emph{fast} and \emph{deadly} mean the same thing in every genre.\\
World Claims & \textcolor{atrAccent2}{\bfseries Workable} & A keep becomes an orbital station. Big, memorable, slow to re-establish.\\
\end{tabularx}
\end{center}

\begin{dread}
The book's own premise is doing design work here. ``What you own is a
\emph{kind} of thing, and you decide what it actually is each time you
enter a new realm'' is already written into the Stuff chapter --- and the
systems that lean hardest into it (Load, Signature Items, Tags) are also
the lightest. The heaviest option, Money, is the one that fights the
premise.
\end{dread}

\section{Vehicles, Armour, and Bases}

These three resist a single treatment because they do different jobs.
Whatever the core Stuff system, each needs an answer:

\begin{opttable}{The Three Awkward Cases}
\tintrow\opt{Vehicles}{Roll the vehicle's rating in place of a Stat; or treat a chase as a clock filled by Piloting rolls; or give the vehicle its own harm track and let passengers hide behind it.}{Cortex, Blades clocks, Traveller}{3}
\opt{Armour}{Reduce incoming harm by rating; or negate one hit per scene; or step harm down one level on a ladder; or add dice to a resistance roll.}{Blades harm ladder, Shadowrun soak}{2}
\tintrow\opt{Bases}{No roll at all. A base grants downtime actions, safety, storage, and a place the story can return to.}{Blades lair, Ars Magica covenant}{2}
\end{opttable}

% =====================================================================
\renewcommand{\realmaccent}{columnSkills}
\chapter{Skills}
% =====================================================================

\lettrine{T}{he current draft} lists seventeen named skills in five
categories. It has a problem worth naming plainly: \textbf{a fixed skill
list carries genre baggage.} \emph{Hacking} is meaningless under a
fantasy sun; \emph{Archery} is dead weight on a space station. Several
options below exist specifically to dodge this.

\begin{opttable}{Seven Ways to Be Good at Things}
\tintrow\opt{Points + Cap}{A budget of points spread over a fixed skill list, capped per skill. \emph{Cap = rank + 1} keeps it to one rule. This is the current direction.}{Vampire: the Masquerade}{3}
\opt{Fixed Array}{Each rank grants a fixed shape --- rank \rating{3} = one at 4, two at 3, three at 1. Assign names to slots. No arithmetic.}{Fate pyramid}{1}
\tintrow\opt{Approaches}{Replace skills with six adverbs: \emph{Careful, Clever, Flashy, Forceful, Quick, Sneaky}. You always have all six; the question is \emph{how} you do a thing. \textbf{Completely genre-proof.}}{Fate Accelerated}{1}
\opt{Freeform Tags}{Invent your own skills in your own words. Rank sets how many and how high. \emph{Same engine the Magic chapter already uses.}}{Lady Blackbird, Risus, Over the Edge}{2}
\tintrow\opt{Careers}{Rate broad life-paths instead of narrow skills --- \emph{Soldier 3, Smuggler 2, Court Wizard 1}. Genre-flexible; a Soldier is a soldier in every realm.}{Traveller, Barbarians of Lemuria, 13th Age}{2}
\opt{Broad + Speciality}{A wide skill at one rating with a narrow speciality that adds dice in its slice. Depth without a longer list.}{Vampire specialities, Shadowrun}{3}
\tintrow\opt{Binary}{You either have a skill or you do not. Having it lets you roll your Stat; lacking it means rolling at a penalty. Stats carry all the maths.}{Into the Odd, Cairn}{1}
\end{opttable}

\section{Breadth Versus Depth}

The Skills chapter already promises this choice at rank \rating{3}. Only
some systems can actually deliver it:

\begin{center}
\setlength{\tabcolsep}{6pt}\renewcommand{\arraystretch}{1.3}
\begin{tabularx}{\textwidth}{@{}p{3.6cm} >{\leavevmode\centering\arraybackslash}p{2.9cm} Y@{}}
\rowcolor{\realmaccent}
\thc{System} &
\thc{Offers the choice?} &
\thc{Why}\\
Points + Cap & \textcolor{columnSkills}{\bfseries Yes} & Spend wide or spend deep. This is precisely what a point budget is for.\\
\rowcolor{atrBoxbg}
Fixed Array & \textcolor{columnBody}{\bfseries No} & The shape is dictated. Fast, but the promised decision disappears.\\
Approaches & \textcolor{columnBody}{\bfseries No} & Everyone has all six. There is no breadth to buy.\\
\rowcolor{atrBoxbg}
Freeform Tags & \textcolor{columnSkills}{\bfseries Yes} & Choose many narrow tags or few broad ones --- and broad tags apply more often.\\
Careers & \textcolor{atrAccent2}{\bfseries Partly} & Careers are inherently broad. Depth comes from stacking related ones.\\
\rowcolor{atrBoxbg}
Broad + Speciality & \textcolor{columnSkills}{\bfseries Yes} & It is the explicit axis of the system.\\
\end{tabularx}
\end{center}

\begin{dread}
Note the coincidence: the Magic chapter already runs on player-invented
tags and credits \emph{Lady Blackbird} for it. If Skills ran on tags too,
the game would have \textbf{one character-building idea applied four
times} rather than four different ones --- and the seventeen-skill list,
with all its genre baggage, would simply evaporate.

The counter-argument is real and should not be waved away: \textbf{a
concrete list is far easier for a new player than a blank line.} Freeform
systems reward experienced players and stall beginners. A hybrid --- a
suggested list that players may ignore --- is available and is what
\emph{Lady Blackbird} itself does.
\end{dread}

% =====================================================================
\renewcommand{\realmaccent}{columnMagic}
\chapter{Magic}
% =====================================================================

\lettrine{M}{agic} is the most finished column: tags chosen at creation,
two dice per applicable tag, a cap on tags per spell set by rank. It is
also the only column that does not use the Stat + Skill frame.

\begin{opttable}{Six Ways to Work Magic}
\tintrow\opt{Tags $\times$ 2}{The current draft. Two dice per applicable tag, up to a rank-set limit. Self-contained, distinctive, already written.}{Lady Blackbird + Ars Magica}{2}
\opt{Tags + Stat}{Pool = tags applied + a Stat. Consistent with every other roll in the game; makes Body matter to casters.}{Mage, Cortex}{2}
\tintrow\opt{Tags as Permission}{Tags do not add dice --- they say what you \emph{can} attempt. Roll Stat + an Arcana skill like anything else. Magic becomes a licence, not an engine.}{Blades, Mage: the Awakening}{1}
\opt{Technique $\times$ Form Grid}{Rate the five Techniques and ten Forms separately; a spell's pool is the sum of the two used. Enormously expressive, genuinely heavy.}{Ars Magica proper}{5}
\tintrow\opt{Always Works, Always Costs}{The spell succeeds. The roll determines what it costs you --- fatigue, corruption, attention, backlash.}{Mage paradox, Shadowrun drain, Ars fatigue}{3}
\opt{Merge with Stuff}{If Stuff also runs on tags, Magic and Stuff become one subsystem wearing two hats: a flaming sword and a fire spell are the same mechanism.}{City of Mist}{2}
\end{opttable}

\section{The Ceiling Problem}

Worth flagging regardless of which option is chosen. Under the current
draft, with one A available per character:

\begin{center}
\setlength{\tabcolsep}{6pt}\renewcommand{\arraystretch}{1.3}
\begin{tabularx}{\textwidth}{@{}p{5.2cm} >{\leavevmode\centering\arraybackslash}p{2.2cm} Y@{}}
\rowcolor{\realmaccent}
\thc{Best available build} &
\thc{Max pool} &
\thc{Note}\\
Body B + Skills A & \dice{9} & Stat 4 + Skill 5.\\
\rowcolor{atrBoxbg}
Stuff A + Skills B & \dice{9} & Only if gear ratings reach 5.\\
Magic A & \dice{10} & Five tags $\times$ two dice, on every spell where five tags apply.\\
\end{tabularx}
\end{center}

Magic reaches a higher ceiling than any mundane build, is not gated
behind owning a specific object, and applies to a far wider range of
problems. Whether that is a bug or the intended meaning of \emph{``you are
a fucking wizard''} is a design decision, not a maths error --- but it
should be a decision.

% =====================================================================
\renewcommand{\realmaccent}{atrAccent}
\chapter{Harm, Armour, and Dying}
% =====================================================================

\lettrine{N}{othing} is decided here at all. The working idea in the book
is ``each rank of armour reduces damage by one,'' with Strength serving as
durability.

\begin{opttable}{Six Ways to Get Hurt}
\tintrow\opt{Hit Points}{A numeric track equal to a Stat or rank. Count down, fall over at zero. Familiar, trivially understood, and utterly genre-neutral.}{D\&D and descendants}{1}
\opt{Harm Ladder}{Four named severities --- Lesser, Moderate, Severe, Fatal. Armour steps a harm down one rung. \textbf{Genre-proof: ``Severe'' is a broken leg in fantasy and a punctured lung in space.}}{Blades in the Dark}{2}
\tintrow\opt{Conditions}{Harm is a fictional state, not a number: \emph{Shaken}, \emph{Wounded}, \emph{Broken}. Each imposes a specific penalty. No arithmetic.}{Apocalypse World, Monster of the Week}{2}
\opt{Stress + Consequences}{A small buffer you can spend to shrug off hits, plus lasting consequences when it runs out. Player-facing choice about when to take a real wound.}{Fate}{3}
\tintrow\opt{Margin as Damage}{Successes above the target \emph{are} the damage. No separate damage roll --- the attack roll does both jobs.}{Shadowrun, Exalted}{2}
\opt{Clocks}{Serious threats are progress tracks the group fills. Death is not a number going to zero; it is a clock completing.}{Blades in the Dark}{2}
\end{opttable}

\begin{opttable}{Four Ways Armour Can Work}
\tintrow\opt{Flat Reduction}{Subtract armour rating from incoming harm. The working idea. Simple, but high armour can trivialise low-harm attacks entirely.}{D\&D damage reduction}{1}
\opt{Step It Down}{Armour reduces harm one level on the ladder, a limited number of times per scene or session. Creates a real ``my armour is used up'' moment.}{Blades in the Dark}{2}
\tintrow\opt{Resistance Dice}{Roll armour rating as a pool to soak. Uncertain, dramatic, and one more roll per hit.}{Shadowrun soak, Vampire}{3}
\opt{Permission}{Armour does not reduce harm; it decides whether a given threat can hurt you at all. Mundane blades simply cannot pierce Master Armour.}{Into the Odd, Nobilis}{2}
\end{opttable}

\begin{dread}
Harm is the one subsystem where the multi-genre premise bites hardest.
A hit-point total implies a consistent scale of injury across realms ---
but a sword, a shotgun, and a plasma lance are not on the same scale
unless they are all just ``harm.'' The \textbf{ladder} and
\textbf{conditions} options sidestep this by never numbering the wound in
the first place.
\end{dread}

% =====================================================================
\renewcommand{\realmaccent}{atrAccent2}
\chapter{Six Complete Packages}
% =====================================================================

\lettrine{T}{he options} above can be mixed freely, but some combinations
are far more coherent than others. Here are six whole games, each
internally consistent, each with a different centre of gravity.

\begin{center}
\setlength{\tabcolsep}{4pt}\renewcommand{\arraystretch}{1.35}
\noindent\begin{tabularx}{\textwidth}{@{}p{1.9cm} Y Y Y Y >{\leavevmode\centering\arraybackslash}p{1.3cm}@{}}
\rowcolor{\realmaccent}
\thc{Package} &
\thc{Core roll} &
\thc{Gear} &
\thc{Skills} &
\thc{Harm} &
\thc{Weight}\\
\optname{A. The Storyteller} & Count successes, 0--10 dice & Substitution & Points + cap, fixed list & Strength as HP & \crunch{4}\\
\rowcolor{atrBoxbg}
\optname{B. The Tag Engine} & Count successes, 2 dice per tag & Tags & Freeform tags & Conditions & \crunch{2}\\
\optname{C. The Forged Realm} & Highest die, 0--4 dice & Load + permission & Careers or actions & Harm ladder & \crunch{2}\\
\rowcolor{atrBoxbg}
\optname{D. The Priority Array} & Count successes, 0--10 dice & N picks & Fixed array & Body rank as HP & \crunch{1}\\
\optname{E. The Toolbox} & Step dice, sum two & Asset die & Points, die sizes & Stress track & \crunch{3}\\
\rowcolor{atrBoxbg}
\optname{F. The Two-Die Realm} & 2d6 + modifier & Modifier & Approaches & Conditions & \crunch{1}\\
\end{tabularx}
\end{center}

\section{A. The Storyteller}
\src{Vampire: the Masquerade, Shadowrun}

The current draft, finished honestly. Four Stats and a skill list, both
built from dot budgets with caps of rank\,+\,1. Gear ratings substitute
for Stats. Magic keeps its tags. Pools run \dice{0}--\dice{10} and the GM
names a target every roll.

\textbf{Best at:} character texture. Two rank-\rating{3} characters can be
completely different people. Optimisers have something to chew on, and
advancement is obvious --- buy another dot.

\textbf{Costs you:} it is the heaviest option, the GM invents a number on
every roll, and the four columns need careful balancing to stay
comparable. The seventeen-skill list keeps its genre baggage.

\section{B. The Tag Engine}
\src{Lady Blackbird, City of Mist}

\textbf{Everything is tags.} Body gives you physical tags, Stuff gives you
possession tags, Skills give you learned tags, Magic gives you arcane
tags. You roll two dice per applicable tag, capped by the rank of
whichever column supplied them. Rank sets how many tags you know and how
many may apply at once.

\textbf{Best at:} radical unification. The whole game is one paragraph.
Genre baggage disappears entirely because players write their own tags.
The Magic chapter is already written in this language, so it comes free.

\textbf{Costs you:} enormous GM adjudication load --- every roll becomes a
negotiation about which tags apply. Blank lines intimidate new players.
And the four columns stop feeling meaningfully different from one another,
which undercuts the 4$\times$4 grid the game is built on.

\section{C. The Forged Realm}
\src{Blades in the Dark}

Rescale everything to \dice{0}--\dice{4}. Read the highest die: 6 is a
full success, 4--5 succeeds with a complication, 1--3 fails. No target
numbers, ever. Gear is a Load you spend during play, not a list you write
beforehand. Harm is a four-rung ladder and armour steps it down.

\textbf{Best at:} speed and the realm-crossing premise. Nothing on the
sheet needs translating between genres because almost nothing on the sheet
is a specific object. Partial success on most rolls keeps stories moving.

\textbf{Costs you:} the \dice{0}--\dice{10} range and everything built on
it, including the current Magic chapter's arithmetic. Rank \rating{1}
versus rank \rating{4} becomes a much smaller gap, which may make the
priority grid feel less consequential.

\section{D. The Priority Array}
\src{D\&D standard array, Fate, Lancer}

Keep the current engine but delete all arithmetic from character creation.
Each rank in each column hands you a \emph{fixed shape}: Body \rating{3}
gives you [4,3,2,1] to arrange; Skills \rating{3} gives you one skill at 4,
two at 3, three at 1; Stuff \rating{3} lets you pick three items rated 3 or
less. Nothing is added up.

\textbf{Best at:} getting to the table. A character takes ten minutes and
cannot be built wrong. Balance across the four columns is guaranteed
rather than hoped for, because you specified the outputs directly.

\textbf{Costs you:} optimisation, and with it some replay depth. Two
Body-\rating{3} characters have identical numbers in different slots.
Advancement needs its own separate design, since there are no points to
spend.

\section{E. The Toolbox}
\src{Cortex Prime, Savage Worlds}

Every trait is a \emph{die size} rather than a count. Body \rating{3}
might grant a d10 and three d6s. Roll your Stat die plus your Skill die
and sum. Gear is simply another die you add or swap in --- which makes the
original question (``a number for the car plus a number for your skill'')
trivially natural, because the car is just a d8.

\textbf{Best at:} gear substitution, which is the question that started
this document. Swapping a d6 for a d10 is intuitive at the table, and
assets, complications, and conditions all slot into the same frame.

\textbf{Costs you:} a rewrite of everything currently expressed in d6
pools, and a subtler probability landscape --- players cannot eyeball odds
the way they can with a coin-flip pool.

\section{F. The Two-Die Realm}
\src{Apocalypse World, Traveller}

Always roll 2d6 and add one small modifier. 6 or less is a miss, 7--9 is a
partial success, 10 or more is a full success. Ranks map to modifiers from
$-1$ to $+3$. Skills become six Approaches; gear adds a modifier or grants
permission.

\textbf{Best at:} absolute minimum friction. Every roll is identical, the
odds are memorised after one session, and partial success is built into
the engine rather than bolted on.

\textbf{Costs you:} the dice pool, which is a real part of the game's
identity --- gathering a fistful of d6s is a pleasure that $2$d$6$ does not
provide. The narrow modifier range also compresses the difference between
rank \rating{1} and rank \rating{4} severely.

% =====================================================================
\renewcommand{\realmaccent}{atrAccent}
\chapter{How to Decide}
% =====================================================================

\section{Dependency Order}

These questions are not independent. Answering them out of order means
redoing work.

\begin{center}
\begin{tikzpicture}[
  node distance=6mm,
  box/.style={draw=atrRule, fill=atrBoxbg, rounded corners=2pt,
              text width=3.1cm, align=center, minimum height=1.05cm,
              font=\small},
  hd/.style={draw=none, fill=atrAccent, text=white, rounded corners=2pt,
              text width=3.1cm, align=center, minimum height=1.05cm,
              font=\hdfont\small},
  arr/.style={-latex, draw=atrRule, line width=0.9pt}]

  \node[hd] (roll) {\textbf{1. Core roll}\\\scriptsize sets the dice range};
  \node[hd, right=12mm of roll] (gear) {\textbf{2. Gear entry}\\\scriptsize sets what Stuff is for};
  \node[box, below left=12mm and 6mm of gear] (body) {\textbf{3a. Body}};
  \node[box, right=4mm of body] (stuff) {\textbf{3b. Stuff}};
  \node[box, right=4mm of stuff] (skills) {\textbf{3c. Skills}};
  \node[box, right=4mm of skills] (magic) {\textbf{3d. Magic}};
  \node[hd, below=12mm of stuff, xshift=14mm] (harm) {\textbf{4. Harm}\\\scriptsize needs 1 and 3a};

  \draw[arr] (roll) -- (gear);
  \draw[arr] (gear.south) -- ++(0,-4mm) -| (body.north);
  \draw[arr] (gear.south) -- ++(0,-4mm) -| (stuff.north);
  \draw[arr] (gear.south) -- ++(0,-4mm) -| (skills.north);
  \draw[arr] (gear.south) -- ++(0,-4mm) -| (magic.north);
  \draw[arr] (body.south) |- ([yshift=6mm]harm.north) -- (harm.north);
  \draw[arr] (stuff.south) |- ([yshift=6mm]harm.north);
\end{tikzpicture}
\end{center}

Chapters 4 through 7 --- the four columns --- are genuinely independent of
one another and can be decided in any order, or by different people.

\section{Five Tests for Any Candidate}

Whatever combination is chosen, run it against these before writing it
into the book.

\begin{statblock}[The Tests]
\begin{tabularx}{\linewidth}{@{}p{3.6cm} Y@{}}
\keyword{The Realm Test} & Take a finished character from high fantasy to the sprawl to a space station. How many lines on the sheet need rewriting? Ideally zero.\\[4pt]
\keyword{The Table Test} & Can a player who has never seen the game build a character in ten minutes, without the GM's help?\\[4pt]
\keyword{The GM Test} & How many numbers must the GM invent per roll? Every one of them is friction, and friction is what makes a game stop being rules-light.\\[4pt]
\keyword{The Column Test} & Play a scene four times, once with each column as the character's A. Does each version feel powerful in a \emph{different} way? If two columns feel the same, one of them is spurious.\\[4pt]
\keyword{The Ceiling Test} & Compute the best possible pool for each column at priority A. Are the four numbers close? If Magic reaches 10 and Stuff reaches 6, Stuff is a trap option.\\
\end{tabularx}
\end{statblock}

\begin{dread}[The one genuinely structural choice]
Most of this document is interchangeable parts. One decision is not.

\textbf{Either the four columns run on four different mechanisms, or they
run on one.} The current draft is already halfway to ``one'' --- Magic's
tag system would work verbatim for Stuff and for Skills. Going all the
way gives a game that fits on an index card, at the cost of the four
columns feeling distinct from each other. Staying at four mechanisms gives
each column its own texture, at the cost of being four times as much game
to learn.

That choice is upstream of nearly everything else here, and it is a
question about what the game is \emph{for} rather than about mathematics.
\end{dread}

\end{document}
